% Part: computability
% Chapter: computability-theory
% Section: reducibility

\documentclass[../../../include/open-logic-section]{subfiles}

\begin{document}

\olfileid{cmp}{thy}{red}
\olsection{راکمېدنه}

\begin{explain}
اوس پوهېږو چې لږ تر لږه يو سټ، $K_0$، شته چې په محاسبوي ډول د شمېر
وړ دے خو محاسبه کېدونکے نۀ دے۔ بايد څرګنده وي چې نور هم شته۔ د
راکمېدنې طريقه يوه پياوړې لار راکوي چې د نورو سټونو دا خاصيتونه، بې
له دې چې هر ځل بنسټيزو اصولو ته ستانۀ شو، ثابت کړو۔

په عمومي ډول، له سټ~$A$ څخه سټ~$B$ ته يوه ``راکمونه'' داسې طريقه ده
چې د دې پوښتنې ځوابونه چې !!{element}s په~$B$ کښې دي که نه، د دې
پوښتنې په ځوابونو بدلوي چې !!{element}s په~$A$ کښې دي که نه۔ موږ به
د ``ډېر-پر-يو راکمېدنې'' په نوم مفهوم ته پام وکړو، خو د بېلابېلو
خاصيتونو لرونکي د راکمېدنې ډېر نور مفهومونه هم شته۔ د راکمېدنې
مفهومونه د محاسبوي پېچلتيا په څېړنه کښې هم مرکزي دي، هلته د اغېزمنتيا
مسئلې هم بايد په پام کښې ونيول شي۔ د بېلګې په توګه، يو سټ هغه وخت
``اېن‌پي-بشپړ'' بلل کېږي چې په اېن‌پي کښې وي او د اېن‌پي هره مسئله
ورته راکمول کېدے شي؛ دلته د لاندې مفهوم ته ورته راکمونه کارېږي، خو
دا زيات شرط لري چې راکمونه په څوحدي وخت کښې محاسبه شي۔

موږ د راکمونې مفهوم لا له مخکښې په ضمني ډول کارولے دے۔ سټ~$K$ داسې
تعريف کړئ:
\[
K = \Setabs{x}{\cfind{x}(x) \fdefined},
\]
يعنې $K = \Setabs{x}{x \in W_x}$۔ زموږ هغه ثبوت چې د درېدنې مسئله
نۀ حل کېږي (\olref[hlt]{thm:halting-problem})، تر ټولو نېغ دا ښيي چې
$K$ محاسبه کېدونکے نۀ دے۔ راياد کړئ چې~$K_0$ دا سټ دے:
\[
K_0 = \Setabs{\tuple{e, x}}{\cfind{e}(x) \fdefined },
\]
يعنې $K_0 = \Setabs{\tuple{e,x}}{x \in W_e}$۔ د~$K$ د نۀ محاسبه
کېدنې هر ثبوت د~$K_0$ نۀ محاسبه کېدنې ته غځول اسان دي: که~$K_0$
محاسبه کېدونکے وای، نو دا مو پرېکړه کولے شوه چې !!a{element}~$x$ په
$K$ کښې دے که نه، يوازې په دې پوښتنې سره چې جوړه~$\tuple{x, x}$ په
$K_0$ کښې ده که نه۔ هغه تابع~$f$ چې~$x$ په~$\tuple{x, x}$ بدلوي،
له~$K$ څخه~$K_0$ ته د \emph{راکمونې} يوه بېلګه ده۔
\textbf{د ژباړې د نسخې سمون (OLCMP-030):} سرچينې د دې بحث په پيل
کښې د ``راکمونې مفهوم'' نوم دوه ځله پرله‌پسې ليکلے او د درېدنې مسئله
ئې د ``نۀ حل کېدونکې'' پر ځاے د يوې ليکدوديزې تېروتنې له امله بله
کلمه بللې وه۔ په پښتو کښې دواړه څرګندې تحريرې تېروتنې نۀ دي نقل شوې۔
\textbf{د ژباړې د نسخې سمون (OLCMP-031):} د درېدنې د جوړو دويم
سټ-جوړونکي بيان په سرچينه کښې د جوړې دوه مختصات اړولي وو، خو شرط ئې
بيا هم دويم عدد د لومړي عدد د شاخص په توګه کاراوه۔ دلته جوړه د لومړي،
کره تعريف په ترتيب ساتل شوې ده؛ ورپسې قطري راکمونه هم همدغه ترتيب
کاروي۔
\end{explain}

\begin{defn}
فرض کړئ $A$ او~$B$ د طبيعي عددونو سټونه دي۔ محاسبه کېدونکې تابع
~$f\colon \Nat \to \Nat$ له~$A$ څخه~$B$ ته \emph{ډېر-پر-يو راکمونه}
ده، هله او يوازې هله چې د هر طبيعي عدد~$x$ دپاره
\[
x \in A \quad \text{هله او يوازې هله چې} \quad f(x) \in B.
\]
که داسې راکمونه~$f$ شته، وايو~$A$، $B$ ته \emph{ډېر-پر-يو راکمېدونکے}
دے، چې $A \leq_m B$ ليکل کېږي۔ که~$A$، $B$ ته ډېر-پر-يو راکمېدونکے
وي او برعکس هم، نو~$A$ او~$B$ \emph{ډېر-پر-يو معادل} بلل کېږي، چې
$A \equiv_m B$ ليکل کېږي۔
\end{defn}

که د پورته تعريف تابع~$f$ اتفاقاً يو په يو وي، نو~$A$، $B$ ته
\emph{يو-پر-يو راکمېدونکے} بلل کېږي۔ لاندې ډېرې راکمونې دا پياوړے
شرط پوره کوي، خو موږ به دا حقيقت نۀ کاروو۔

\begin{digress}
دا رښتيا ده، خو هېڅ څرګنده نۀ ده، چې يو-پر-يو راکمېدنه په رښتيا د
ډېر-پر-يو راکمېدنې په پرتله پياوړے شرط دے۔ په بله وينا، داسې نامتناهي
سټونه~$A$ او~$B$ شته چې~$A$، $B$ ته ډېر-پر-يو راکمېدونکے وي خو
~$B$ ته يو-پر-يو راکمېدونکے نۀ وي۔
\end{digress}

\end{document}
